% ==============================================================================
\chapter{Conclusiones y trabajo futuro}
\label{ch:conclusiones}
% ==============================================================================

\section{Logros}
\label{sec:logros}

Repasando los ocho objetivos planteados en la sección~\ref{sec:objetivos}, siete de ellos están plenamente cumplidos: la arquitectura, el modelo de datos, la API REST documentada, la interfaz PWA, los mecanismos de seguridad del Sprint~6, el análisis del estado del arte y la documentación reproducible. El objetivo~7 (validación con pirámide de pruebas y usuarios reales) queda parcialmente cumplido, ya que la pirámide de pruebas automatizadas (unitarias, integración y E2E) está operativa, pero la validación de usabilidad con usuarios reales sigue pendiente y se desarrolla como línea de trabajo futuro.



\subsubsection*{Transparencia con el uso de IA}
  El proyecto es coherente con el contexto en el que se desarrolla. Existen herramientas de asistencia y las he usado, declarándolo y dejando trazas auditables. La sección~\ref{sec:uso-ia} y el Capítulo~\ref{ch:implementacion} entran en detalle. Creo que esa actitud es la única razonable a día de hoy.

\subsubsection*{Calidad cercana a producción}
  El sistema incluye CI/CD con GitHub Actions, \textit{audit log} de seguridad, documentación API autogenerada, capacidad \textit{offline} y 2FA. Estas funcionalidades dan lugar a un resultado más profesional, aunque no eran estrictamente necesarias para cumplir el enunciado.

\subsubsection*{Reusabilidad}
  El modelo de datos, el flujo de estados y el patrón de delegación a entidades pueden trasladarse con pocos cambios a dominios cercanos: seguridad ciudadana, patrimonio, salud pública. El diseño no se queda en lo medioambiental.


\section{Limitaciones}
\label{sec:limitaciones}

Reconocer lo que no está bien terminado es parte del trabajo:

\subsubsection*{Cobertura de pruebas}
  La cobertura de \textit{branches} (38{,}91\,\%) está claramente por debajo del 70\,\% que fijé en RNF-08. Muchas ramas de error no se prueban. Es un trabajo mecánico al que no se le ha dado tiempo para cerrar.

\subsubsection*{Sin SMTP real}
  La recuperación de contraseña por email (RF-05) y las notificaciones por email (RF-46) están programadas pero no se han probado contra un servidor real. Funcionarán al configurar las variables de entorno en producción.

\subsubsection*{Escalado por \textit{timeout} no programado}
  El RF-34 (escalado de incidencias no atendidas) está diseñado pero no hay un \textit{scheduler} corriendo. Falta añadir cron o un \textit{worker} dedicado.

\subsubsection*{Pruebas de usabilidad pendientes}
  La validación con usuarios reales sigue el protocolo descrito en la sección~\ref{sec:usabilidad}, pero su ejecución completa ha quedado fuera del alcance de esta entrega. Una evaluación sólida requeriría entre 5 y 15 participantes con perfiles diversos y se aborda como trabajo futuro inmediato.

\subsubsection*{Sin prueba de carga}
  El RNF-02 (100 usuarios simultáneos) se ha validado por dimensionado teórico, no con \textit{k6} o \textit{Artillery} contra el sistema real.

\subsubsection*{Moderación social básica}
  Los comentarios entre usuarios funcionan, pero no hay filtros de lenguaje, reputación ni \textit{shadow banning}. En un despliegue público abierto haría falta.


\section{Trabajo futuro}
\label{sec:trabajo-futuro}

Las líneas que veo claras como continuación natural:

\textbf{Corto plazo.}  Subir la cobertura por encima del 80\,\%, desplegar en \textit{cloud} con HTTPS real (Let's Encrypt) y backups automatizados, programar el \textit{scheduler} de RF-34, y pasar pruebas de carga con \textit{k6} para validar RNF-02.

\textbf{Medio plazo.}  Una aplicación nativa complementaria con React Native, aprovechando la API REST existente, para mejorar las notificaciones \textit{push} y el acceso a sensores. Asignación automática de incidencias a entidades por intersección geoespacial entre ubicación y jurisdicción (el campo \texttt{jurisdiction} ya existe en el modelo). Sistema de reputación del ciudadano basado en histórico. Clasificación automática de la fotografía mediante visión por computador, que sugiera categoría y severidad al crear la incidencia.

\textbf{Largo plazo.}  Convenios con ayuntamientos piloto en la provincia de Alicante para validar el producto en producción. Integración con plataformas de gestión municipal (GESCON, Gestiona). Modelo de negocio B2G \textit{freemium} desarrollado a partir del análisis del Anexo~\ref{app:modelo-negocio}. Apertura del código con licencia MIT o Apache~2.0 para fomentar replicación en otros municipios.


\section{Reflexión}
\label{sec:reflexion}

El TFG ha permitido juntar en un único producto cosas que vimos en distintas asignaturas del Grado: bases de datos, programación web, arquitectura, ingeniería de requisitos, pruebas, seguridad, gestión de proyectos. Hasta ahora cada una vivía en su silo; aquí se han tenido que coordinar.

En lo metodológico, combinar SDD con agentes de IA en un IDE como Antigravity ha sido un aprendizaje en sí mismo. Lo más probable es que sea así como va a trabajar buena parte de la industria de aquí a unos años. La conclusión que se obtiene es que la diferencia no la marca delegar, sino dirigir bien: especificar con precisión, revisar con rigor y firmar el resultado.

Por último, se agradece al tutor José Luis Sánchez Romero su disponibilidad y la libertad que me dio para elegir el \textit{stack} y el alcance. Sin esa libertad, EcoAlerta no habría salido como ha salido.
